\PassOptionsToPackage{no-math}{fontspec}
\documentclass[UTF8,zihao=-4,openany,twoside,fontset=fandol]{ctexbook}
\usepackage[paperwidth=190mm,paperheight=270mm,inner=24mm,outer=21mm,top=25mm,bottom=25mm,headheight=16pt,headsep=8mm,marginparwidth=15mm,marginparsep=3mm]{geometry}
\usepackage{fontspec}
\setmainfont{texgyrepagella}[Extension=.otf,UprightFont=*-regular,BoldFont=*-bold,ItalicFont=*-italic,BoldItalicFont=*-bolditalic]
\setsansfont{texgyreheros}[Extension=.otf,UprightFont=*-regular,BoldFont=*-bold,ItalicFont=*-italic,BoldItalicFont=*-bolditalic,Scale=MatchLowercase]
\usepackage{amsmath,amssymb,amsthm,mathtools,newpxmath}
\SetSymbolFont{operators}{normal}{OT1}{zpllf}{m}{n}
\SetSymbolFont{operators}{bold}{OT1}{zpllf}{b}{n}
\SetMathAlphabet{\mathbf}{normal}{OT1}{zpllf}{b}{n}
\SetMathAlphabet{\mathit}{normal}{OT1}{zpllf}{m}{it}
\SetMathAlphabet{\mathit}{bold}{OT1}{zpllf}{b}{it}
\usepackage{graphicx,xcolor,tikz,tikz-cd,booktabs,longtable,array,enumitem,microtype,fancyhdr,needspace,etoolbox,multicol}
\usetikzlibrary{arrows.meta,calc,positioning,decorations.pathmorphing,patterns}
\definecolor{Ink}{HTML}{143B49}
\definecolor{Gold}{HTML}{B38C4F}
\definecolor{Muted}{HTML}{687D82}
\definecolor{Mist}{HTML}{DCE8E5}
\definecolor{Paper}{HTML}{F4F1E9}
\usepackage[unicode,colorlinks,linkcolor=Ink,citecolor=Ink,urlcolor=Ink,bookmarksnumbered,pdfdisplaydoctitle,hypertexnames=false]{hyperref}
\hypersetup{pdftitle={代数几何基础（FGA）中文译本},pdfauthor={Alexander Grothendieck},pdfsubject={1957—1962年Bourbaki报告；中英法术语双向索引}}
\urlstyle{same}
\ctexset{chapter={format=\LARGE\bfseries\color{Ink},name={第,篇},number=\chinese{chapter},beforeskip=5pt,afterskip=22pt},section={format=\Large\bfseries\color{Ink},beforeskip=19pt,afterskip=9pt},subsection={format=\large\bfseries\color{Ink},beforeskip=14pt,afterskip=7pt},subsubsection={format=\normalsize\bfseries\color{Ink}}}
\setcounter{secnumdepth}{0}
\setcounter{tocdepth}{1}
\setlength{\parindent}{2em}
\setlength{\parskip}{2pt}
\linespread{1.18}
\raggedbottom
\setlength{\emergencystretch}{3em}
\allowdisplaybreaks[2]
\setlist[itemize]{leftmargin=2em,itemsep=3pt,topsep=5pt}
\setlist[enumerate]{leftmargin=2.6em,itemsep=3pt,topsep=5pt}
\setlist[description]{style=nextline,font=\normalfont\bfseries,leftmargin=1.8em,itemsep=5pt}
\newif\iftermanchors\termanchorstrue
\pagestyle{fancy}\fancyhf{}
\fancyhead[LE]{\small\color{Muted}代数几何基础\quad FGA}
\fancyhead[RO]{\small\color{Muted}\begingroup\termanchorsfalse\nouppercase{\leftmark}\endgroup}
\fancyfoot[LE,RO]{\color{Ink}\thepage}
\fancyfoot[RE,LO]{\scriptsize\color{Muted}A. Grothendieck}
\renewcommand{\headrulewidth}{0.3pt}
\renewcommand{\chaptermark}[1]{\markboth{\ifcase\value{chapter}#1\or 凝聚代数层的对偶\or 形式几何与代数几何\or 下降 I：忠实平坦下降\or 下降 II：形式模理论\or 下降 III：商预概形\or 下降 IV：Hilbert 概形\or 下降 V：Picard 概形的存在\or 下降 VI：Picard 概形的性质\else #1\fi}{}}
\fancypagestyle{plain}{\fancyhf{}\fancyfoot[LE,RO]{\color{Ink}\thepage}\renewcommand{\headrulewidth}{0pt}}
\newcommand{\chapterintro}[2]{\vspace{-9pt}\begin{flushleft}\color{Muted}\small\itshape #1\par\normalfont\small #2\end{flushleft}\vspace{12pt}}
\newcommand{\origpage}[2]{\par\phantomsection\label{orig:#1}\marginpar{\raggedright\scriptsize\color{Muted}原书\par #2}}
\makeatletter
\newcommand{\frenchterm}[1]{{\color{Muted}\fontsize{\dimexpr\f@size pt*9/10\relax}{\baselineskip}\selectfont（#1）}}
\makeatother
\DeclareRobustCommand{\term}[3]{#1\frenchterm{#3}}
\DeclareRobustCommand{\termuse}[4]{\iftermanchors\leavevmode\allowbreak\mbox{\hypertarget{termocc-#2}{}\label{termocc:#2}\hyperlink{term-#1}{#3}}\frenchterm{#4}\else #3\fi}
\newenvironment{statement}[1]{\par\addvspace{9pt}\needspace{3\baselineskip}\noindent\textbf{\color{Ink}#1}\quad\ignorespaces}{\par\addvspace{7pt}}
\renewcommand{\proofname}{证明}
\newcounter{translatornote}
\DeclareRobustCommand{\translatornote}[1]{\footnote{\textbf{译者注。}#1}}
\DeclareRobustCommand{\translatornoteuse}[2]{\stepcounter{translatornote}\footnote{\hypertarget{tn-#1}{}\label{tn:#1}\textbf{\hyperlink{tn-index-#1}{译者注\arabic{translatornote}。}}#2}}
\newcommand{\authorcorrection}[1]{\footnote{\textbf{作者勘误。}#1}}
\newcommand{\Spec}{\operatorname{Spec}}
\newcommand{\Spf}{\operatorname{Spf}}
\newcommand{\Hom}{\operatorname{Hom}}
\newcommand{\Ext}{\operatorname{Ext}}
\newcommand{\Tor}{\operatorname{Tor}}
\newcommand{\Pic}{\operatorname{Pic}}
\newcommand{\Hilb}{\operatorname{Hilb}}
\newcommand{\Gal}{\operatorname{Gal}}
\newcommand{\Aut}{\operatorname{Aut}}
\newcommand{\Isom}{\operatorname{Isom}}
\newcommand{\End}{\operatorname{End}}
\newcommand{\Ker}{\operatorname{Ker}}
\newcommand{\Coker}{\operatorname{Coker}}
\newcommand{\coker}{\operatorname{Coker}}
\newcommand{\im}{\operatorname{Im}}
\newcommand{\rank}{\operatorname{rank}}
\newcommand{\Tr}{\operatorname{Tr}}
\newcommand{\id}{\operatorname{id}}
\newcommand{\NS}{\operatorname{NS}}
\newcommand{\Alb}{\operatorname{Alb}}
\newcommand{\Supp}{\operatorname{Supp}}
\newcommand{\Proj}{\operatorname{Proj}}
\newcommand{\Sym}{\operatorname{Sym}}
\newcommand{\depth}{\operatorname{depth}}
\newcommand{\Char}{\operatorname{char}}
\pdfstringdefDisableCommands{\def\term#1#2#3{#1}\def\termuse#1#2#3#4{#3}\def\translatornote#1{}\def\translatornoteuse#1#2{}\def\mathcal#1{#1}\def\mathbb#1{#1}\def\Pic{Pic}\def\Hilb{Hilb}\def\Spec{Spec}}
\begin{document}
\frontmatter
\begin{titlepage}
\thispagestyle{empty}
\begin{tikzpicture}[remember picture,overlay]
\fill[Ink] (current page.south west) rectangle (current page.north east);
\begin{scope}[shift={(current page.south west)},x=1cm,y=1cm]
\fill[white,opacity=.025] (0,10) rectangle (19,18);
\draw[Gold,line width=.7pt] (2,24.2)--(17,24.2);
\draw[Mist,opacity=.32,line width=.5pt] (2,8.4)--(17,8.4);
\foreach \j in {0,...,6}{
 \pgfmathsetmacro{\xx}{2.7+2.15*\j}
 \draw[Mist,opacity=.18,line width=.35pt] (\xx,8.4)--(\xx,16.3);
 \begin{scope}[shift={(\xx,11.5)},xscale=.65,yscale=.7]
  \draw[Gold,opacity=.65,line width=.6pt] (-1.15,-1.25) .. controls (-.3,-.2) and (-.6,.6) .. (0,1.35) .. controls (.6,.6) and (.3,-.2) .. (1.15,-1.25);
  \draw[Mist,opacity=.50,line width=.5pt] (-1.15,.65) .. controls (-.5,.2) and (-.55,-.4) .. (0,-.5) .. controls (.55,-.4) and (.5,.2) .. (1.15,.65);
 \end{scope}
 \fill[Gold] (\xx,8.4) circle (1.15pt);
}
\node[anchor=west,text=Paper,font=\fontsize{11}{14}\selectfont\sffamily] at (2,25.2) {ALEXANDER GROTHENDIECK};
\node[anchor=north west,text=Paper,font=\fontsize{41}{53}\selectfont\bfseries,align=left] at (1.8,22.9) {代数几何\\基础};
\node[anchor=west,text=Gold,font=\fontsize{12}{17}\selectfont\itshape] at (2,17.8) {Fondements de la géométrie algébrique};
\node[anchor=west,text=Mist,font=\small] at (2,16.7) {Bourbaki 报告选集\quad 1957—1962};
\node[anchor=west,text=Paper,font=\fontsize{34}{40}\selectfont\sffamily] at (2,6.9) {FGA};
\node[anchor=west,text=Mist,font=\small] at (2,5.7) {中文译本\quad 附中·英·法术语索引};
\node[anchor=east,text=Mist,font=\small] at (17,2.4) {对偶\enspace ·\enspace 下降\enspace ·\enspace 模空间};
\draw[Gold,line width=.7pt] (2,3.35)--(17,3.35);
\end{scope}
\end{tikzpicture}
\null\vfill
\end{titlepage}
\cleardoublepage
\thispagestyle{empty}
\null\vfill
\begin{flushleft}
{\Large\bfseries\color{Ink}代数几何基础}\par\medskip
Alexander Grothendieck 著\par
\textit{Fondements de la géométrie algébrique}\par
\textit{Extraits du Séminaire Bourbaki 1957--1962}\par\bigskip
原版：Secrétariat mathématique，11 rue Pierre Curie，Paris 5e，1962。\par
底本：用户提供的 208 页扫描本。\par
排版：XeLaTeX；正文与数学公式重新排印，交换图以 TikZ 重绘。\par
本译本为辅助阅读的中文翻译。\par
\end{flushleft}
\vfill
\clearpage
\chapter*{译本说明}
\addcontentsline{toc}{chapter}{译本说明}
本书汇集 Grothendieck 于1957至1962年在 Bourbaki 讨论班所作的八篇报告，以及1962年补写的评论与勘误。译本依底本次序收录全部内容，保留各篇的节号、命题号和公式编号。页边的“原书”标记指向法文印刷页码，例如“190-07”表示第190号报告的第7页；扫描本中149-04有破损页和完整重扫页两张，译文采用完整页，不重复收录。

技术术语在中文后括注法文。点击正文中的术语，可跳至书末中、英、法对照表；表中的各个出现位置可分别跳回正文。为方便和其他文献对读，历史用语 \textit{préschéma} 译作“预概形”，\textit{schéma} 译作“概形”；理解其含义时仍应以本书各处的定义和假设为准。旧称 \textit{simple sur} 的相关用法，按上下文译作“在……上光滑”，与“正则”区别处理。

作者的评论反映1962年时的进展，其中“最近”“目前”“尚未发表”等表述均依原文保留。作者已指出的错误或缺失条件，在相应正文处说明其依据，并链接到书前评论。另加的数学说明以“译者注”标明；这类说明给出补充条件、具体反例或论证，不与原文混排为作者的结论。书末另列译者注索引，便于集中查阅。

本书的报告带有提纲和研究预告的性质；原文未给出的证明，译文不自行写成原作者已完成的证明。文献条目保留原题名，同时补出中文题意，便于检索原文。

\begingroup\termanchorsfalse\tableofcontents\endgroup
\input{generated/00-commentary}
\mainmatter
\input{generated/01-duality}
\input{generated/02-formal}
\input{generated/03-descent}
\input{generated/04-existence}
\input{generated/05-quotient}
\input{generated/06-hilbert}
\input{generated/07-picard-existence}
\input{generated/08-picard-properties}
\backmatter
\renewcommand{\chaptermark}[1]{\markboth{#1}{}}
\input{generated/translator-notes}
\input{generated/glossary}
\end{document}
